% Capítulo textual do trabalho.

\chapter{Códigos e Algoritmos}
\label{cap:codigos}

Este capítulo apresenta o uso de algoritmos e trechos de código, abordando desde a inclusão básica até a personalização de estilos e numeração. A formatação clara e precisa de algoritmos e códigos é fundamental em áreas como ciência da computação, matemática e engenharias, onde a transparência dos métodos e a reprodutibilidade são essenciais para a validade dos resultados apresentados~\cite{knuth1984texbook}.

\section{Inclusão de Códigos}

Para incluir códigos de programação, o pacote mais comum é o \texttt{listings}. Esse pacote suporta várias linguagens de programação e permite a personalização do estilo de exibição, incluindo o realce de sintaxe, a numeração de linhas e a adição de comentários. O ambiente \texttt{codigo} formata a posição do código e localização da legenda e da fonte.

O Código~\ref{cod:codigo} apresenta um exemplo de código para inclusão de códigos-fonte e o Código~\ref{cod:fatorial} mostra o código renderizado.
Este código implementa uma função recursiva em Python para calcular o fatorial de um número. A sintaxe da linguagem Python é destacada automaticamente e o código é numerado.

\begin{codigo}[htb]
\legenda[cod:codigo]{Código para inclusão de códigos-fonte}
\begin{lstlisting}
\begin{codigo}[htb]
    \legenda[cod:codigo]{Código para inclusão de códigos-fonte}
    \begin{lstlisting}[language=Python]
    def fatorial(n):
        if n == 0:
            return 1
        else:
            return n * fatorial(n - 1)
    \end{lstlisting }
    \fonteautor
\end{codigo}
\end{lstlisting}
\fonteautor
\end{codigo}

\begin{codigo}[htb]
\legenda[cod:fatorial]{Função recursiva para calcular o fatorial de um número}
\begin{lstlisting}[language=Python]
def fatorial(n):
    if n == 0:
        return 1
    else:
        return n * fatorial(n - 1)
\end{lstlisting}
\fonteautor
\end{codigo}

\section{Inclusão de Algoritmos}

A formatação de algoritmos é realizada de maneira análoga à inclusão de códigos, porém, utiliza o ambiente \texttt{algoritmo}. O Algoritmo~\ref{alg:fatorial} mostra um exemplo de algoritmo para o cálculo do fatorial de maneira recursiva.

\begin{algoritmo}
\legenda[alg:fatorial]{Algoritmo para o calcular o fatorial de um número}
\begin{lstlisting}[mathescape]
ENTRADA: Número inteiro não negativo $n$
SAÍDA Valor de $n!$

SE $n = 0$
    RETORNA $1$
SENÃO
    RETURNA $n \times \text{fatorial}(n-1)$
FIM-SE
\end{lstlisting}
\fonteautor
\end{algoritmo}

\section{Considerações Finais}

O uso de algoritmos e códigos é um recurso poderoso para a formatação de trabalhos acadêmicos e científicos, especialmente em áreas técnicas. A personalização oferecida pelo LaTeX permite a adaptação dos estilos de exibição conforme as necessidades específicas do autor ou da instituição.
